\chapter{Knee Pain}

\section*{Definition}
Pain localized to the knee joint and surrounding structures, one of the most common musculoskeletal complaints affecting mobility and quality of life.

\section*{Pathophysiology}
The knee's weight-bearing role and complex ligamentous/cartilage structures make it vulnerable to acute injury (ligament tears, meniscal tears) and chronic degenerative changes (osteoarthritis). Anterior knee pain suggests patellofemoral pathology; medial/lateral pain suggests meniscal or collateral ligament involvement; diffuse pain suggests inflammatory arthritis.

\section*{Causes}
\begin{itemize}
  \item Osteoarthritis (degenerative joint disease)
  \item Meniscal tear
  \item Anterior cruciate ligament (ACL) injury
  \item Patellofemoral pain syndrome
  \item Patellar tendinopathy (jumper's knee)
  \item Baker's cyst (popliteal cyst)
  \item Rheumatoid arthritis or gout
  \item Iliotibial band syndrome
  \item Prepatellar bursitis (housemaid's knee)
  \item Septic arthritis or osteomyelitis
\end{itemize}

\section*{When to See a Doctor}
See your doctor if:
\begin{itemize}
  \item Severe or worsening symptoms
  \item Acute knee injury with inability to bear weight, significant swelling, fever, locking, or instability
  \item Accompanied by other concerning signs
\end{itemize}

\section*{Related Symptoms}
\begin{itemize}
  \item Swelling
  \item Stiffness
  \item Locking or catching
  \item Instability (giving way)
  \item Limping
\end{itemize}
\textbf{Important:} If this symptom persists, your doctor will check for serious underlying causes.

\section*{Self-Care / Home Management}
\begin{itemize}
  \item Rest and avoid activities that make it worse.
  \item Maintain adequate hydration and a balanced diet.
  \item Over-the-counter remedies may provide symptomatic relief (consult a pharmacist or doctor).
  \item Monitor symptoms and seek professional advice if they persist or worsen.
  \item Avoid self-medicating with prescription medications without medical guidance.
\end{itemize}

\section*{How Your Doctor Will Evaluate You}
\begin{itemize}
  \item A thorough history and physical examination are the first steps in evaluation.
  \item Laboratory tests (blood work, urinalysis) may be ordered based on clinical suspicion.
  \item Imaging studies (X-ray, ultrasound, CT, MRI) may be indicated when structural pathology is suspected.
  \item Specialist referral is arranged when the diagnosis remains uncertain or requires specific expertise.
\end{itemize}

\section*{Treatment Options}
\begin{itemize}
  \item Treatment targets the underlying cause identified through diagnostic evaluation.
  \item Supportive care and symptom management are provided while awaiting definitive therapy.
  \item Medications, lifestyle modifications, and physical therapies may be prescribed as appropriate.
  \item Follow-up ensures treatment effectiveness and allows timely adjustments to the care plan.
\end{itemize}

\clearpage